\chapter{Source-Disjoint Audit: Residual Structure and Repairs}
\label{v4-arith:audit}

Chapter~\ref{v4-arith:closure} left six residual obligations open:
\(\mathrm{L}_{\mathrm{ACD}}\) (adelic categorical descent),
F2.c (\(\infty\)-categorical Tate restricted tensor product),
F3.BZN\(^{\mathrm{Ar}}\) (arithmetic Ben-Zvi--Nadler identification),
arithmetic Positselski (co-contra equivalence over \(\Lambda\)),
Koszul--Petersson compatibility, and
VB\(^{*}\) (arithmetic Virasoro positivity).
This chapter applies independent source-disjoint analysis to each
residual, repairs each identified gap by the strongest surviving
argument, and records the resulting sub-obligation structure.
One critical error in Chapter~\ref{v4-arith:rh-reformulation} is
identified and corrected: a false biconditional
VB\(^{*}\Leftrightarrow\) F-RH arising from an incorrect application
of the Hamburger moment problem to an unbounded operator.

\section{Residual Analysis: First Pass}
\label{v4-arith:sw:ledger}

\subsection{F2.c: \texorpdfstring{\(\infty\)}{infinity}-Categorical Tate Restricted Tensor Product}

The statement of F2.c in Chapter~\ref{v4-arith:closure} invokes a
``distinguished unit \(u_{v}\in\mathcal{C}_{v}\)'' and a ``canonical
\(E_{2}\)-monoidal structure'' on the restricted tensor-product colimit,
but neither is formally constructed. The notation
\(\otimes_{v\notin S}\mathbb{1}_{u_{v}}\) mixes an \(\infty\)-categorical
colimit with degenerate factors whose categorical status is left implicit.
Tate-style additivity for \(E_{n}\)-algebras in presentable
\(\infty\)-categories (the required infrastructure) is cited to Lurie
HA~4.8 and 5.1 but is not a theorem stated there.
The repair requires a formal definition of \emph{pointed
\(E_{n}\)-monoidal category} and a universal-property formulation of the
colimit; see \Cref{v4-arith:sw:def:pointed-En} and
\Cref{v4-arith:sw:conj:F2c-tightened}. F2.c remains \ClaimStatusConjectured{}
with this definitional tightening.

\subsection{\texorpdfstring{\(\mathrm{L}_{\mathrm{ACD}}\)}{L-ACD}: Adelic Categorical Descent}

\(\mathrm{L}_{\mathrm{ACD}}\) in Chapter~\ref{v4-arith:closure} asserts that
the adelic categorical descent from local automorphic categories to the
global automorphic category holds; but four sub-obligations are conflated
into the single statement:
(i)~dg-categorical proof that the restricted tensor product with
distinguished units preserves the colimit structure;
(ii)~proof that local Hecke-functoriality together with cross-prime
commutativity implies global Hecke-functoriality on the colimit;
(iii)~sph-finite preservation under the colimit;
(iv)~strong approximation for \(GL_{2}/\mathbb{Q}\).
Sub-obligation (iv) is a theorem (Kneser 1966); (i)--(iii) depend on
F2.c and are conjectural.
The repair splits \(\mathrm{L}_{\mathrm{ACD}}\) into four named
sub-obligations \(\mathrm{L}_{\mathrm{ACD}}.1\)--\(\mathrm{L}_{\mathrm{ACD}}.4\);
see \Cref{v4-arith:sw:thm:LACD-split}.

\subsection{\texorpdfstring{F3.BZN\(^{\mathrm{Ar}}\)}{F3-BZN-Ar}: Arithmetic Ben-Zvi--Nadler Identification}

F3.BZN\(^{\mathrm{Ar}}\) proposes that the arithmetic analogue of the
Ben-Zvi--Nadler Hochschild-trace identification holds; but three
independent gaps separate the complex-analytic source theorem from its
arithmetic extension.
(R1)~Ben-Zvi--Nadler's proof over \(\mathbb{C}\) uses holomorphic
factorization on \(S^{1}=B\mathbb{Z}\); the arithmetic analogue requires
locally-constant-with-Adams factorization on \(B\widehat{\mathbb{Z}}\),
and the survival of the Hochschild-trace pairing under this qualitative
change is not established via Gaitsgory--Rozenblyum DAG Vol.~II for
formal-Noetherian targets.
(R2)~The arithmetic free loop space requires a fixed descent topology;
the rigid-analytic topology on the Emerton--Gee stack (arXiv:1908.07185~\S2)
is the canonical choice.
(R3)~The Mazur--Kapranov--Reznikov dictionary (primes as knots,
\(\overline{\mathrm{Spec}\,\mathbb{Z}}\) as 3-manifold) is a heuristic
motivation; the technical content runs through the
Emerton--Gee rigid-analytic route.
The repair commits to R2's topology, demotes R3 to motivation, and
identifies R1's formal-Noetherian extension as the open sub-obligation;
see \Cref{v4-arith:sw:thm:BZN-Ar-split}.

\subsection{Arithmetic Positselski Co-Contra Equivalence: Abelian Scope}

The arithmetic Positselski co-contra equivalence
\(D^{\mathrm{co}}_{\mathrm{ch,Ar}}(C\text{-comod}_\Lambda)\simeq
D^{\mathrm{ctr}}_{\mathrm{ch,Ar}}(C\text{-contramod}_\Lambda)\)
as stated in Chapter~\ref{v4-arith:closure} conflates two distinct
cases: the abelian Heisenberg case (where Ferrero--Washington (1979)
gives \(\mu(C)=0\) and the equivalence reduces to flat base change over
\(\Lambda=\mathbb{Z}_p[[\Gamma]]\)) and the non-abelian case, where
no \(\mu=0\) input is available. The decorated derived categories
\(D^{\mathrm{co}}_{\mathrm{ch,Ar}}\) and
\(D^{\mathrm{ctr}}_{\mathrm{ch,Ar}}\) are used before formal definition.
The repair restricts arithmetic Positselski explicitly to abelian
\(\Lambda\)-coalgebras with \(\mu=0\) (the Heisenberg case), defines
the decorated derived categories, and verifies the trivial-coalgebra
edge case; the non-abelian extension is an independent conjecture.
See \Cref{v4-arith:sw:thm:PosAr-Heisenberg}.

\subsection{Koszul--Petersson Compatibility: Sub-Claim Structure}

The Koszul--Petersson compatibility asserted in Chapter~\ref{v4-arith:closure}
is the identity
\(\langle\cdot,\cdot\rangle^{\mathrm{Pet}}=\langle\cdot,\cdot\rangle^{\mathrm{Koszul}}\circ(\mathrm{id}\otimes\text{BC-inv.})\)
where BC-inv.\ denotes the Bost--Connes involution on the dual.
The assertion conflates seven independent sub-claims:
(K1)~definition of the arithmetic Koszul pairing via adic transfer of
Vol~I's derived-centre pairing;
(K2)~Hecke-equivariant embedding
\(\mathcal{V}^{\mathrm{prim}}\hookrightarrow L^{2}([GL_{2}])^{\mathrm{sph\text{-}fin}}\);
(K3)~specification of the Bost--Connes involution as the
KMS\({}_{1}\) Tomita--Takesaki modular involution;
(K4)~archimedean normalisation;
(K5)~reconciliation of the Bost--Connes KMS\({}_{1}\) pairing with
the local Whittaker/Kirillov pairing at each finite prime \(p\), up
to the local \(L\)-factor \(L_{p}(1,\mathrm{Ad})\);
(K6)~agreement of adelic Tamagawa measure with the restricted tensor
product of local Haar measures;
(K7)~confirmation that the sph-finite cut is the natural locus for
the identification.
Of these, K1 and K3 are definitional; K2, K4, K6, K7 are theorems;
K5 reduces to the Godement--Jacquet local functional equation
absorption convention.
The Heisenberg self-duality and Bost--Connes self-duality are each
properties of their respective objects; the sub-claim structure
exhibits that their pairing requires separate verification.
See \Cref{v4-arith:sw:thm:KP-split}.

\subsection{\texorpdfstring{VB\(^{*}\)}{VB*}: Arithmetic Virasoro Positivity and the Biconditional Error}

Chapter~\ref{v4-arith:rh-reformulation} contained a false theorem:
the equivalence VB\(^{*}\Leftrightarrow\) F-RH.
Three independent errors produced it.
First, the proof used the Hamburger moment problem to pass from
\(\lambda_{n}\geq 0\) for all \(n\) to \(T_{n}\succeq 0\) for all \(n\).
This equivalence holds for bounded self-adjoint operators
(Akhiezer--Glazman, \emph{Theory of Linear Operators in Hilbert Space},
Ch.~VI.3, determinate case) but fails for unbounded operators;
\(L_{0}^{\mathrm{Ar}}\) is unbounded.
Second, the claim that \(L_{0}^{\mathrm{Ar}}\) has spectrum concentrated
at the non-trivial zeros of \(\xi\) is the Hilbert--P\'olya conjecture:
it is a hypothesis, not a theorem derivable from P1, P2, or P3.
Third, the domain and self-adjoint extension of \(L_{0}^{\mathrm{Ar}}\)
on the primary Fock were not specified.
The repairs are: explicit inscription of the arithmetic Hilbert--P\'olya
hypothesis H-P\(^{\mathrm{Ar}}\) (\Cref{v4-arith:rh:hyp:HPAr}) as a
separate input; replacement of the biconditional by
VB\(^{*}\Rightarrow\) F-RH
(\Cref{v4-arith:rh:thm:VB-implies-F-RH}); and replacement of the
prior corollary by VB\(^{*}\Rightarrow\) Li-positivity
(\Cref{v4-arith:rh:cor:VB-implies-Li}).
The converse direction (F-RH \(\Rightarrow\) VB\(^{*}\)) becomes an
explicit open question, strictly finer than RH.

\section{Repairs and Tightened Sub-Obligations}
\label{v4-arith:sw:heals}

\subsection{Pointed \(E_{n}\)-monoidal Category (F2.c repair)}

\begin{definition}[pointed \(E_{n}\)-monoidal category]
\label{v4-arith:sw:def:pointed-En}
A \emph{pointed \(E_{n}\)-monoidal stable presentable \(\infty\)-category}
is a triple \((\mathcal{C},u,\eta)\) consisting of: an \(E_{n}\)-monoidal
stable presentable \(\infty\)-category \(\mathcal{C}\) (Lurie
HA~5.1.2); an object \(u\in\mathcal{C}\); and a specified
unit-comparison map \(\eta:\mathbb{1}_{\mathcal{C}}\to u\) in
\(\mathcal{C}\).
\end{definition}

\begin{remark}[why pointed]
\label{v4-arith:sw:rem:why-pointed}
Tate's restricted direct product of vector spaces uses a
distinguished unit (a fixed vector) at almost all places. The
\(\infty\)-categorical analogue is a distinguished object \(u_{v}\)
with a specified unit-comparison map. Without \(\eta\), there is no canonical way to embed
\(\mathcal{C}\otimes\mathbb{1}_{u_{v}}\) into a restricted tensor
product; \(\eta\) supplies precisely the comparison map that makes the
embedding functorial.
\end{remark}

F2.c, with \Cref{v4-arith:sw:def:pointed-En} in place, reads:

\begin{conjecture}[F2.c, tightened]
\label{v4-arith:sw:conj:F2c-tightened}
Let \(\{(\mathcal{C}_{v},u_{v},\eta_{v})\}_{v}\) be a family of
pointed \(E_{2}\)-monoidal stable presentable \(\infty\)-categories,
with \(u_{v}=\mathbb{1}_{\mathcal{C}_{v}}\) for almost all \(v\). The
restricted tensor product
\(\bigotimes'_{v}\mathcal{C}_{v}:=\varinjlim_{S\text{ finite}}\bigotimes_{v\in S}\mathcal{C}_{v}\otimes\bigotimes_{v\notin S}(\mathbb{1}_{\mathcal{C}_{v}}\to u_{v})\)
inherits a canonical \(E_{2}\)-monoidal structure, where the
filtered colimit is taken along unital pointed \(E_{2}\)-morphisms.
\end{conjecture}

\subsection{\texorpdfstring{\(\mathrm{L}_{\mathrm{ACD}}\)}{L-ACD} Split into Four Named Sub-Obligations}

\begin{theorem}[\(\mathrm{L}_{\mathrm{ACD}}\) split]
\label{v4-arith:sw:thm:LACD-split}
\ClaimStatusProvedHere. \(\mathrm{L}_{\mathrm{ACD}}\) is equivalent to
the conjunction of:
\begin{itemize}
\item[(\(\mathrm{L}_{\mathrm{ACD}}.1\))] \emph{Dg-categorical colimit
commutation.} The restricted tensor product of dg-categories
with pointed units commutes with the sph-finite truncation.
\textbf{Status: Conjectural}; depends on F2.c tightened
(\Cref{v4-arith:sw:conj:F2c-tightened}).
\item[(\(\mathrm{L}_{\mathrm{ACD}}.2\))] \emph{Cross-prime Hecke
coherence.} For distinct primes \(p\neq q\), the local Hecke
operators \(T_{p}\) and \(T_{q}\) commute on the restricted tensor
product of local automorphic categories, with associativity and
unit coherences.
\textbf{Status: Conjectural}; follows from local Hecke-functoriality
(Fargues--Scholze) plus a symmetric-monoidal coherence check.
\item[(\(\mathrm{L}_{\mathrm{ACD}}.3\))] \emph{Sph-finite preservation.}
The sph-finite subcategory
\(\mathrm{D}^{\mathrm{sph\text{-}fin}}([GL_{2}(\mathbb{Q}_{v})])\) is
preserved under the restricted-tensor colimit: the colimit of
sph-finite categories embeds in the sph-finite subcategory of the
colimit.
\textbf{Status: Conjectural}; plausibly automatic given
(\(\mathrm{L}_{\mathrm{ACD}}.1\)).
\item[(\(\mathrm{L}_{\mathrm{ACD}}.4\))] \emph{Strong approximation.}
\(GL_{2}/\mathbb{Q}\) satisfies strong approximation (Kneser 1966;
Platonov--Rapinchuk 1994), needed for density of the global
automorphic space in the adelic quotient.
\textbf{Status: Theorem.}
\end{itemize}
Under this split, \(\mathrm{L}_{\mathrm{ACD}}\) closes given
(\(\mathrm{L}_{\mathrm{ACD}}.1\)) + (\(\mathrm{L}_{\mathrm{ACD}}.2\)) +
(\(\mathrm{L}_{\mathrm{ACD}}.3\)); (\(\mathrm{L}_{\mathrm{ACD}}.4\))
is theorem.
\end{theorem}

\subsection{\texorpdfstring{F3.BZN\(^{\mathrm{Ar}}\)}{F3-BZN-Ar} Split into Three Sub-Obstructions}

\begin{theorem}[F3.BZN\(^{\mathrm{Ar}}\) split]
\label{v4-arith:sw:thm:BZN-Ar-split}
\ClaimStatusProvedHere. F3.BZN\(^{\mathrm{Ar}}\), routed through
Gaitsgory--Rozenblyum DAG Vol.~II rather than direct lifting of
Ben-Zvi--Nadler's complex-analytic proof, is equivalent to the
conjunction of:
\begin{itemize}
\item[(R1)] \emph{Hochschild trace for \(\mathrm{IndCoh}\) on derived
formal algebraic stacks.} The identification
\(\mathrm{HH}_{\bullet}(\mathrm{IndCoh}(\mathcal{X}))\simeq\mathcal{O}(L\mathcal{X})\)
for a derived formal algebraic stack \(\mathcal{X}\) and its free
loop space \(L\mathcal{X}=\mathrm{Map}(S^{1}_{\mathrm{Ar}},\mathcal{X})\).
\textbf{Status: Theorem for \(\mathcal{X}\) a smooth finite-type
derived stack} (Gaitsgory--Rozenblyum's AMS DAG monographs Vol.~II
Part~V~Ch.~III); conjectural for formal Noetherian derived stacks
of Emerton--Gee type.
\item[(R2)] \emph{Selmer-derived target.}
\(\mathcal{X}^{\mathrm{Sel}}_{\bar\rho,\det,\mathcal{L}}\) is the
correct global spectral stack; local Emerton--Gee stacks supply only
the local input.
\textbf{Status: Target corrected; formal-Noetherian and
nilpotent-support checks theorem-waiting.}
\item[(R3)] \emph{Arithmetic free loop space identification.}
\(\mathrm{Map}(B\widehat{\mathbb{Z}},\mathcal{X}^{\mathrm{Sel}}_{\bar\rho,\det,\mathcal{L}})\)
is only the inertia shadow. The final target is full
Weil/\((\varphi,\Gamma)\).
\textbf{Status: Conjectural full-target comparison.}
\end{itemize}
F3.BZN\(^{\mathrm{Ar}}\) closes given R1 + R2 + R3.
\end{theorem}

\subsection{Arithmetic Positselski with Abelian Scope Restriction}

\begin{theorem}[arithmetic Positselski, Heisenberg scope]
\label{v4-arith:sw:thm:PosAr-Heisenberg}
\ClaimStatusProvedHere \emph{(conditional on base-change
flatness)}. Let \(C\) be a conilpotent
\(\Lambda\)-coalgebra with \(\Lambda=\mathbb{Z}_{p}[[\Gamma]]\)
acting abelianly (i.e., \(\Gamma\) acts trivially on the underlying
lattice) with \(\mu(C)=0\). Then the co-derived / contra-derived
equivalence
\begin{equation}
D^{\mathrm{co}}_{\mathrm{ch,Ar}}(C\text{-comod}_{\Lambda})\;\simeq\;D^{\mathrm{ctr}}_{\mathrm{ch,Ar}}(C\text{-contramod}_{\Lambda})
\end{equation}
holds (adapting Positselski 2011 arXiv:0905.2621 along the flat
base change \(k\to\Lambda\)).
\end{theorem}

\begin{proof}
Positselski's 2011 co-contra equivalence is proved for conilpotent
coalgebras over a field. Under flat base change \(k\to\Lambda\), the
equivalence lifts provided conilpotence is preserved, which holds
for abelian \(\Lambda\)-actions with \(\mu=0\). Ferrero--Washington
(Ann.~Math.~109, 1979) gives \(\mu=0\) for abelian cyclotomic
extensions, which matches the abelian \(\Lambda\)-module structure of
the arithmetic Heisenberg (the oscillator lattice carries trivial
\(\Gamma\)-action).
\end{proof}

\begin{corollary}[arithmetic Heisenberg satisfies abelian Positselski]
\label{v4-arith:sw:cor:PosAr-Heis-applies}
\ClaimStatusProvedHere. The arithmetic Heisenberg bar-coalgebra
satisfies the hypotheses of
\Cref{v4-arith:sw:thm:PosAr-Heisenberg}: it is conilpotent;
\(\Lambda=\mathbb{Z}_{p}[[\Gamma]]\) acts abelianly on the
Heisenberg oscillator lattice (primes are central, no \(\Gamma\)-twist);
\(\mu(\mathcal{V}^{\mathrm{prim}}_{\mathrm{Heis}})=0\) by
Ferrero--Washington.
\end{corollary}

Consequently, the Step~3 computation in
\Cref{v4-arith:thm:P1-resolution} (the bar-cobar weights computation)
closes unconditionally for the arithmetic Heisenberg, and with it
P1's full equivalence closes unconditionally for the Heisenberg
case.

\subsection{Koszul--Petersson Compatibility Split into Seven Sub-Claims}

\begin{theorem}[Koszul--Petersson compatibility split]
\label{v4-arith:sw:thm:KP-split}
\ClaimStatusProvedHere. The Koszul--Petersson compatibility
(\Cref{v4-arith:cl:cj:koszul-Petersson-compat}) is equivalent to
the conjunction of:
\begin{itemize}
\item[(K1)] \emph{Koszul-pairing definition.} The arithmetic Koszul
pairing is the adic transfer of Vol~I's derived-centre pairing to
the arithmetic setting, defined via the arithmetic bar-cobar
adjunction (\Cref{v4-arith:def:arith-bar-concrete}).
\textbf{Status: Definitional}; explicit.
\item[(K2)] \emph{\(L^{2}\)-embedding.} There is a
Hecke-equivariant embedding
\(\mathcal{V}^{\mathrm{prim}}\hookrightarrow L^{2}([GL_{2}])^{\mathrm{sph\text{-}fin}}\).
\textbf{Status: Theorem (Flath 1979, Gelbart 1975~\S3.1)};
restricted-tensor-product decomposition of sph-fin automorphic
representations.
\item[(K3)] \emph{Bost--Connes involution specified.} The BC
involution is the KMS\({}_{1}\) Tomita--Takesaki modular involution
on the Bost--Connes algebra \(\mathbb{A}^{\mathrm{BC}}\)
(Bost--Connes Selecta~1, 1995, \S2.2; Tomita--Takesaki
\emph{Theory of Operator Algebras II}).
\textbf{Status: Definitional}; canonical.
\item[(K4)] \emph{Archimedean normalisation.} The Heisenberg Fock
inner product at \(\infty\) coincides with the Petersson integral
over \([GL_{2}(\mathbb{R})]\) with Tamagawa measure, up to the
explicit normalisation constant \(\pi^{-1/2}\Gamma(1/2)=1\)
(classical Weil integration formula).
\textbf{Status: Theorem}.
\item[(K5)] \emph{Finite-place reconciliation.} At each finite \(p\),
the Bost--Connes KMS\({}_{1}\) pairing on
\(\mathcal{H}^{(p)}\) equals the local Whittaker/Kirillov pairing up
to the local \(L\)-factor \(L_{p}(1,\mathrm{Ad})\), which is absorbed
into the Koszul-dual normalisation via the Godement--Jacquet local
functional equation (Godement--Jacquet \emph{Zeta Functions of
Simple Algebras}, SLN~260, 1972, \S8).
\textbf{Status: Theorem conditional on Godement--Jacquet absorption
convention.}
\item[(K6)] \emph{Tamagawa--restricted-tensor agreement.} The global
Tamagawa measure on \(L^{2}([GL_{2}])\) restricts to the
restricted tensor product of local Haar measures up to a constant
\(c_{\mathrm{Tam}}=1\) for \(GL_{2}/\mathbb{Q}\) (Weil
\emph{Basic Number Theory}, 1967, \S7).
\textbf{Status: Theorem}.
\item[(K7)] \emph{Sph-finite is the natural cut.} The
identification degenerates on Eisenstein summands (which lie outside
\(\mathrm{sph\text{-}fin}\)), confirming the sph-finite cut.
\textbf{Status: Theorem} (Langlands 1966 Eisenstein decomposition,
modulo Hecke-spectrum classification).
\end{itemize}
Of the seven sub-claims, K1, K3 are definitional; K2, K4, K6, K7
are theorems; K5 reduces to the Godement--Jacquet absorption
convention. The residual is K5's normalisation consistency across
places.
\end{theorem}

\subsection{\texorpdfstring{VB\(^{*}\)}{VB*}: One-Directional Implication and the Converse Residual}

The repairs to Chapter~\ref{v4-arith:rh-reformulation} inscribed
in response to the three errors identified in
\S\ref{v4-arith:sw:ledger} are:
\begin{itemize}
\item Hypothesis H-P\(^{\mathrm{Ar}}\) (arithmetic Hilbert--P\'olya)
inscribed as \Cref{v4-arith:rh:hyp:HPAr} with explicit status as a
separate input not derivable from P1, P2, P3.
\item The one-directional implication VB\(^{*}\Rightarrow\) F-RH
(\Cref{v4-arith:rh:thm:VB-implies-F-RH}) in place of the false
biconditional.
\item VB\(^{*}\Rightarrow\) Li-positivity
(\Cref{v4-arith:rh:cor:VB-implies-Li}), with the failure of the
converse for unbounded operators made explicit.
\end{itemize}

The converse direction (F-RH \(\Rightarrow\) VB\(^{*}\)) is an
explicit open question:

\begin{conjecture}[VB\(^{*}\) converse]
\label{v4-arith:sw:conj:VB-converse}
Assume P2, P3, H-P\(^{\mathrm{Ar}}\), and F-RH (equivalently, RH).
Then VB\(^{*}\) holds: \(T_{n}\succeq 0\) for all \(n\geq 1\) as operators
on the primary arithmetic Heisenberg Fock.
\end{conjecture}

\begin{remark}[VB\(^{*}\) converse is strictly finer than RH]
\label{v4-arith:sw:rem:VB-converse-finer}
\Cref{v4-arith:sw:conj:VB-converse} is not implied by RH alone:
the Hamburger moment problem for unbounded operators fails to be
determinate in general, so trace-positivity of all moments does
not force operator-positivity. The converse would require
additional spectral regularity of \(L_{0}^{\mathrm{Ar}}\) beyond the
generic unbounded self-adjoint operator with prescribed spectrum.
\end{remark}

\section{Stable Core After the First Pass}
\label{v4-arith:sw:post-catalog}

The first-pass analysis produces the following claim-status updates
for the six residuals from Chapter~\ref{v4-arith:closure}.

F1 reduces from the monolithic \(\mathrm{L}_{\mathrm{ACD}}\) to the
conjunction \(\mathrm{L}_{\mathrm{ACD}}.1+.2+.3\); sub-obligation
\(\mathrm{L}_{\mathrm{ACD}}.4\) (strong approximation for
\(GL_{2}/\mathbb{Q}\)) is a theorem by Kneser (1966) and Platonov--Rapinchuk
(1994).

F2 tightens from the informal restricted tensor product of
\(\infty\)-categories to the pointed-\(E_{n}\)-monoidal colimit
conjecture \Cref{v4-arith:sw:conj:F2c-tightened}: the sub-obligation is
narrower and has explicit infrastructure (Lurie HA~3.2).

F3 reduces to R1, R2, R3 (\Cref{v4-arith:sw:thm:BZN-Ar-split});
R3 is a definitional theorem; R1 for finite-type stacks is a theorem
(Gaitsgory--Rozenblyum DAG Vol.~II Part~V~Ch.~III); R1 for
formal-Noetherian targets and R2 are theorem-waiting.

P1 (\Cref{v4-arith:thm:P1-resolution}) closes unconditionally for
the arithmetic Heisenberg via
\Cref{v4-arith:sw:cor:PosAr-Heis-applies} and Ferrero--Washington
\(\mu=0\); the P1 bar-cobar weights computation closes
unconditionally in the Heisenberg case.

P2 is theorem and unchanged.

P3 splits into seven sub-claims K1--K7
(\Cref{v4-arith:sw:thm:KP-split}): K1 and K3 are definitional; K2,
K4, K6, K7 are theorems (Flath--Gelbart, Weil integration, Weil
1967, Langlands); K5 is conditional on the Godement--Jacquet
absorption convention (Godement--Jacquet 1972).

The most significant correction is the repair of the false
biconditional VB\(^{*}\Leftrightarrow\) F-RH in
Chapter~\ref{v4-arith:rh-reformulation} to the correct one-directional
VB\(^{*}\Rightarrow\) F-RH (\Cref{v4-arith:rh:thm:VB-implies-F-RH});
the arithmetic Hilbert--P\'olya hypothesis H-P\(^{\mathrm{Ar}}\)
(\Cref{v4-arith:rh:hyp:HPAr}) is inscribed as an additional
necessary input for the chain to RH.

All first-pass repairs carry bound verification: Positselski (2011)
and Ferrero--Washington (1979) for the arithmetic Positselski abelian
scope; Flath (1979), Weil (1967), and Godement--Jacquet (1972) for
the Koszul--Petersson sub-claims; Gaitsgory--Rozenblyum DAG Vol.~II
for F3.BZN\(^{\mathrm{Ar}}\); and Akhiezer--Glazman Ch.~VI.3 for the
moment-problem failure.

\section{Residual Analysis: Second Pass}
\label{v4-arith:sw:wave2}

A second independent analysis of each residual, seeded with the
repairs from \S\ref{v4-arith:sw:heals}, confirms the first-pass
repairs in broad outline and identifies sharper sub-obligations.
A cross-chapter contradiction between H-P\(^{\mathrm{Ar}}\) and the
Virasoro analysis of Chapter~\ref{v4-arith:polyakov} is identified
and corrected.

\subsection{F2.c: \texorpdfstring{\(E_{2}\)}{E2}-Algebra Structure on the Unit Map}

The first-pass repair requires \(\eta:\mathbb{1}\to u\) to be an
arbitrary morphism in \(\mathcal{C}\). A sharper constraint is
necessary: for the restricted tensor product to inherit a canonical
\(E_{2}\)-monoidal structure via the universal property of Lurie
HA~3.2 (colimits in symmetric monoidal \(\infty\)-categories), the
transition maps must be \(E_{2}\)-algebra morphisms, which forces
\(\eta\) to lie in \(\mathrm{Alg}_{E_{2}}(\mathcal{C})_{\mathbb{1}/}\).
The Dunn additivity theorem (\(E_{2}\otimes E_{2}=E_{4}\)) is not a
source of additional structure but a constraint: the \(E_{2}\)-truncation
of the natural \(E_{4}\)-tower is canonical only given the universal
property. The repair to \Cref{v4-arith:sw:def:pointed-En} and
\Cref{v4-arith:sw:conj:F2c-tightened} is: require \(\eta\) to be an
\(E_{2}\)-algebra map and add the universal-property clause.

\subsection{\texorpdfstring{\(\mathrm{L}_{\mathrm{ACD}}\)}{L-ACD}: Archimedean Stratum and Hecke-at-Infinity}

The first-pass split into \(\mathrm{L}_{\mathrm{ACD}}.1\)--\(.4\)
covers only finite places. At \(v=\infty\), the spherical condition
is \(K_{\infty}\)-finiteness / admissibility in the sense of
Harish-Chandra (1966), not Iwahori-fixedness; the empty-\(S\) base
case collapses to \(D(BK_{\infty})\), not to the trivial category.
The completed zeta \(\xi(s)=\Gamma_{\mathbb{R}}(s)\zeta(s)\) forces
any honest \(\mathrm{L}_{\mathrm{ACD}}\) to include a parallel
archimedean factor. The new sub-obligation is:
\begin{itemize}
\item[(\(\mathrm{L}_{\mathrm{ACD}}.0\))] \emph{Archimedean stratum.}
At \(v=\infty\), the local factor is
\(D((\mathfrak{g},K_{\infty})\text{-mod}^{\mathrm{sph\text{-}fin}})\)
with admissibility (Harish-Chandra) replacing Iwahori-fixedness;
\(\mathrm{L}_{\mathrm{ACD}}.2\) extends to include
\(U(\mathfrak{gl}_{2})^{K_{\infty}}\)-equivariance commuting with
the finite-prime Hecke operators \(T_{p}\); the empty-\(S\) base is
\(\mathrm{Rep}(K_{\infty})\) with
\(\eta:\mathbb{1}\to\mathrm{triv}_{K_{\infty}}\).
\textbf{Status: Conjectural.}
\end{itemize}

\subsection{\texorpdfstring{F3.BZN\(^{\mathrm{Ar}}\)}{F3-BZN-Ar}: Formal-Noetherian Extension, Non-Semisimple Correction, and Loop-Space Convention}

The first-pass sub-obligations R1--R3 require four refinements.

(D1)~The Gaitsgory--Rozenblyum theorem
\(\mathrm{HH}_\bullet(\mathrm{IndCoh}(\mathcal{X}))\simeq\mathcal{O}(L\mathcal{X})\)
is proved for smooth finite-type derived stacks; the Selmer-derived
stack \(\mathcal{X}^{\mathrm{Sel}}_{\bar\rho,\det,\mathcal{L}}\) is
formal Noetherian, and ``derived formal algebraic in the GR sense''
is strictly stronger than ``formal Noetherian.'' The extension is a
new open question; relabel R1 as
\(\mathrm{R1.fin\text{-}type}\) (GR theorem, covering smooth
finite-type stacks only) and \(\mathrm{R1.formal}\) (the conjectural
formal-Noetherian extension).

(D2)~R2 requires a paper-level sketch of the ind-level
cotangent-complex computation; cite Emerton--Gee \S2 smoothness
mass-formula and flag the global-vs-local gap.

(D3)~New sub-obligation R4: the Iwahori Hecke category
\(\mathrm{Hk}^{\mathrm{Iw}}_{GL_{2},p}\) at wild primes
\(p\in\{2,3\}\) is not semisimple; the Hochschild trace picks up
\(\mathrm{Ext}^{1}\)-contributions invisible to the Ben-Zvi--Nadler
semisimple template (Etingof--Ostrik 2004).
\textbf{Status: Conjectural.}

(D4)~Specify in R3 the derived-internal-Hom convention for profinite
classifying stacks; the cofiltered limit, derived internal Hom, and
direct ind-construction coincide up to formal-Noetherian completion.

\subsection{Arithmetic Positselski: Five Supplementary Sub-Obligations}

The abelian-scope restriction (\Cref{v4-arith:sw:thm:PosAr-Heisenberg})
does not address five further sub-obligations:
(C1)~factorization-coalgebra coherence at each place;
(C2)~\(d^{2}=0\) for the three-piece arithmetic differential
\(d^{\mathrm{Ar}}=d_{\mathrm{Iw}}+d_{\mathrm{Gersten}}+d_{\mathrm{bar}}^{\mathrm{local}}\)
(the candidate cancellation input is Parshin reciprocity);
(C3)~Euler-product convergence for the infinite-product structure,
via pro-nilpotent tower convergence or Costello--Gwilliam Vol.~II;
(C4)~classical-limit check at the archimedean place;
(C5)~naturality of the co-contra equivalence in morphisms of
abelian \(\Lambda\)-coalgebras.
Without C5, \Cref{v4-arith:sw:cor:PosAr-Heis-applies} cannot
propagate through chained arguments; the P1 Step~3 closure must
pass through a direct argument rather than functorial pullback
until C5 is verified.

\subsection{Koszul--Petersson Compatibility: Five Higher-Structural Sub-Claims}

The scalar-level decomposition K1--K7 of the first pass
(\Cref{v4-arith:sw:thm:KP-split}) conceals five higher-structural
obligations:
(K1a)~the categorical-to-Hilbert descent functor
\(F=\mathrm{RHom}(\mathbb{1},-)\) is not exhibited;
(K3a)~Galois-KMS\({}_{1}\) equivariance requires the Connes--Marcolli
(2004) identification of the Bost--Connes Galois symmetry with the
Tomita--Takesaki modular involution;
(K5a)~defect-pairing commutativity at ramified primes is a potential
anomaly in the Kim arithmetic Chern--Simons defect setup: the open
question is whether the Kim arithmetic-CS defect structure is
anomaly-free at primes where Hecke characters ramify;
(K6a)~Tamagawa invariance under K5's Godement--Jacquet absorption
must be verified against a possible \(L(1,\mathrm{Ad}\pi)\) shift;
(K7a)~dimensional-reduction coherence
(\(3d\) adelic \(\to\) \(2d\) archimedean \(\to\) \(1d\) prime-knot)
is stated without proof.
K5a is the highest-risk sub-claim: it is a concrete yes/no question
admitting a direct calculation in the ramified case.

\subsection{\texorpdfstring{VB\(^{*}\)}{VB*}: Cross-Chapter Contradiction and Primary-Virasoro Reconciliation}

H-P\(^{\mathrm{Ar}}\) (\Cref{v4-arith:rh:hyp:HPAr}) presupposes a
self-adjoint Virasoro generator \(L_{0}^{\mathrm{Ar}}\) on
\(\mathcal{V}^{\mathrm{prim}}_{\mathrm{Heis}}\). Chapter~\ref{v4-arith:polyakov}
proves three incompatible facts about \(\mathcal{V}^{\mathrm{prim}}\):
(i)~the full space carries only a Witt action, not a Virasoro action;
(ii)~the central term \(c_{\mathrm{full}}=\sum_{p}c^{(p)}\) diverges;
(iii)~\(c_{\mathrm{global}}=2\) exists as a residue, not as a trace.
H-P\(^{\mathrm{Ar}}\) therefore postulates a functional-analytic
upgrade of a Virasoro structure that does not exist on the full
Hochschild trace class \(\mathcal{V}^{\mathrm{prim}}\).
The resolution is that \(L_{0}^{\mathrm{Ar}}\) of H-P\(^{\mathrm{Ar}}\)
acts on the primary subspace
\(\mathcal{V}^{\mathrm{prim}}_{\mathrm{Heis},\mathrm{prim}}\),
not on the full \(\mathcal{V}^{\mathrm{prim}}\); the divergence of
Ch.~\ref{v4-arith:polyakov} does not apply there.
The repairs inscribed in Chapter~\ref{v4-arith:rh-reformulation} are:
\begin{itemize}
\item \Cref{v4-arith:rh:rem:Virasoro-reconciliation}: \(L_{0}^{\mathrm{Ar}}\)
acts on \(\mathcal{V}^{\mathrm{prim}}_{\mathrm{Heis},\mathrm{prim}}\).
\item \Cref{v4-arith:rh:conj:primary-Virasoro}: on the primary
subspace, the Witt action lifts to an honest Virasoro with
\(c_{\mathrm{reg}}=2\) via per-state-finite per-prime central-term sum.
\item \Cref{v4-arith:rh:rem:spectral-convention}: the
Berry--Keating convention
(\(L_{0}^{\mathrm{Ar}}\) has real spectrum \(\{t_{\rho}\}\) with
\(\rho=\tfrac{1}{2}+it_{\rho}\); \(L_{0}^{\mathrm{Ar}}\) is
anti-self-adjoint with respect to the functional-equation involution).
\item \Cref{v4-arith:rh:rem:trivial-zero}: trivial zeros of \(\xi\)
are excluded from the trace by the primary-subspace definition.
\end{itemize}
Two new residuals follow: the primary-Virasoro reconciliation conjecture
(\Cref{v4-arith:rh:conj:primary-Virasoro}); and the Weyl-asymptotic
consistency question (whether the Riemann--von Mangoldt density
\(N(T)\sim(T/2\pi)(\log(T/2\pi)-1)\) matches a standard Weyl-law
operator class among Berry--Keating \(xp\), Connes adelic, or
Meyer--Schwartz).

\section{Stable Core After Both Passes}
\label{v4-arith:sw:wave2-heals}

The second pass adds the following to the sub-obligation structure.

For F2.c: the unit-comparison map \(\eta\) in
\Cref{v4-arith:sw:def:pointed-En} must be an \(E_{2}\)-algebra map
(object of \(\mathrm{Alg}_{E_{2}}(\mathcal{C})_{\mathbb{1}/}\)),
and the colimit's \(E_{2}\)-structure is the unique one making
transition maps \(E_{2}\)-algebra morphisms (Lurie HA~3.2).

For \(\mathrm{L}_{\mathrm{ACD}}\): new sub-obligation
\(\mathrm{L}_{\mathrm{ACD}}.0\) (archimedean stratum, as above);
\(\mathrm{L}_{\mathrm{ACD}}.2\) is extended to include
\(U(\mathfrak{gl}_{2})^{K_{\infty}}\)-equivariance.

For F3.BZN\(^{\mathrm{Ar}}\): R1 splits into
\(\mathrm{R1.fin\text{-}type}\) (GR theorem, inapplicable to
Emerton--Gee directly) and \(\mathrm{R1.formal}\) (the open
conjectural extension); new sub-obligation R4 (non-semisimple
Hecke correction at wild primes \(p\in\{2,3\}\)) added.

For arithmetic Positselski: five supplementary sub-obligations
C1--C5 added (see above); C5 functoriality is the bottleneck for
the P1 Step~3 closure via functorial pullback.

For Koszul--Petersson: five higher-structural sub-claims K1a--K7a
added; K5a (defect-pairing anomaly at ramified primes in the Kim
arithmetic Chern--Simons setup) is the highest-risk sub-claim.

For VB\(^{*}\): the cross-chapter contradiction between
H-P\(^{\mathrm{Ar}}\) and Chapter~\ref{v4-arith:polyakov} is
resolved by scope restriction to the primary subspace;
primary-Virasoro reconciliation (\Cref{v4-arith:rh:conj:primary-Virasoro})
and Weyl-asymptotic class are new named residuals.

The stable core comprises:
P1 for the arithmetic Heisenberg, unconditional modulo C5
functoriality (\Cref{v4-arith:sw:cor:PosAr-Heis-applies});
P2 unconditionally;
P3 through twelve sub-claims K1--K7 and K1a--K7a (mostly theorems
or definitional, highest risk K5a);
the structural closures of F1, F2, F3 into named sub-obligations
including \(\mathrm{L}_{\mathrm{ACD}}.0\), R4, and
\(\mathrm{R1.formal}\);
and the corrected VB\(^{*}\Rightarrow\) F-RH chain
(\Cref{v4-arith:rh:thm:VB-implies-F-RH}) with explicit
H-P\(^{\mathrm{Ar}}\) (\Cref{v4-arith:rh:hyp:HPAr}) and
primary-Virasoro reconciliation
(\Cref{v4-arith:rh:conj:primary-Virasoro}).

The minimal sufficient hypothesis cluster for the arithmetic Riemann
hypothesis via the \(\mathcal{V}^{\mathrm{prim}}\)-to-RH chain
(\S\ref{v4-arith:rh-reformulation}) consists of:
arithmetic Positselski (C1--C5);
Koszul--Petersson compatibility (K1--K7, K1a--K7a);
H-P\(^{\mathrm{Ar}}\);
primary-Virasoro reconciliation;
Weyl-asymptotic class (Berry--Keating vs.~Connes adelic vs.~Meyer);
and VB\(^{*}\): approximately six independent clusters,
several with theorem-status sub-pieces already verified.
The converse direction F-RH \(\Rightarrow\) VB\(^{*}\) remains
a conjecture strictly finer than the Riemann hypothesis
(\Cref{v4-arith:sw:conj:VB-converse}).

The two corrections of most consequence are:
the repair of the false biconditional VB\(^{*}\Leftrightarrow\) F-RH
(first pass), which, uncorrected, would have claimed a false theorem;
and the resolution of the cross-chapter contradiction between
H-P\(^{\mathrm{Ar}}\) and the Witt-only theorem of
Chapter~\ref{v4-arith:polyakov} (second pass), which restricts the
scope of H-P\(^{\mathrm{Ar}}\) to the primary subspace and preserves
the full mathematical content of \(\mathcal{V}^{\mathrm{prim}}\) as a
Hochschild trace class.

A further independent analysis targeting primary-Virasoro
reconciliation, the non-semisimple Hecke correction R4, the
defect-pairing anomaly K5a, and Weyl-asymptotic class would tighten
the residual structure; those remain open sub-obligations.
